% ===== 原文 pp.23--25（PDF pp.24--26）=====
% tr-03 片段 a：4.3 末尾「非 Calabi--Yau 情形」、4.4 标题、CP$^{N-1}$ 模型（到 (4.28)）
% 拼装顺序：tr03a → tr03b → tr03c

% [原文 p.23]
\para{非 Calabi--Yau 情形}

设上述条件被破坏，即 $b_1\neq 0$。此时 FI 参数在重整化群下跑动——当能标由 $\mu$ 变到 $\mu'$ 时，它按
\begin{equation}
    \zeta \to \zeta' = \zeta + \log(\mu'/\mu)\,b_1, \tag{4.21}
\end{equation}
变化；而经典的轴 $U(1)$ R-对称性是有异常的——轴向转动 $\ee^{\ii\beta}\in U(1)_A$ 会把 theta 角平移为
\begin{equation}
    \theta \to \theta + 2\beta b_1. \tag{4.22}
\end{equation}
% [Tutor Note: p.23 式 (4.21) 的 $b_1$ 是 FI 参数的一圈 beta 函数系数，等于各物质场规范荷之和 $\sum_i Q_i$（不是 $\sum_i Q_i^2$，后者是规范异常系数）。可验证的锚点：CP$^{N-1}$ 模型的 $N$ 个荷为 1 的场给出 $b_1=N$，正是 p.24 的 "$\zeta'=\zeta+N\log(\mu'/\mu)$"；T$^{U(1)}_{N,d}$ 模型的荷为 $-d,1,\dots,1$ 给出 $b_1=N-d$，正是 p.26 的 "$\zeta'=\zeta+(N-d)\log(\mu'/\mu)$"。(4.22) 中的 $2b_1$ 与 (4.21) 同源：取 $\tW_{\mathrm{eff}}=-t\sigma-\sum_iQ_i\sigma(\log(Q_i\sigma/\mu)-1)$，轴向转动使 $\sigma\to\ee^{2\ii\beta}\sigma$，即 $\log(Q_i\sigma/\mu)\to\log(Q_i\sigma/\mu)+2\ii\beta$，于是 $\tW_{\mathrm{eff}}\to\tW_{\mathrm{eff}}+2\ii\beta b_1\sigma$；与 $\tW_{\mathrm{eff}}=-t\sigma$ 比较得 $\delta t=-2\ii\beta b_1$，再由 $t=\zeta-\ii\theta$ 即得 (4.22)。]
相应地，(4.16) 中的 $\tW_{\mathrm{eff}}$ 会非平凡地依赖 $\mu$。由于真空方程 $\dd\tW_{\mathrm{eff}}(\sigma)\in 2\pi\ii\mathrm{P}$ 没有标度不变性，Coulomb 真空的集合不必是非紧的。相当常见的情形是存在孤立的 Coulomb 真空。这样的真空无法通过对 $U(\sigma,\phi)$ 的经典分析找到，但必须作为量子理论的一个合法真空加以计入。当然，也可以存在 $\sigma$ 取特殊值（例如 $\sigma=0$）的真空，它们可以由经典分析找到。

理论在红外极限下流向某个孤立的 Coulomb 真空（它通常是一个有质量真空），或者流向 $\sigma=0$ 处的 Higgs 分支理论 $(\mu^{-1}(\zeta_{\mathrm{IR}})/G,W_{\zeta_{\mathrm{IR}}})$——其中 $\zeta_{\mathrm{IR}}$ 是 FI 参数的红外取值——或者流向这两类情形的混合。某些 Higgs 分支理论可以是非平凡的 SCFT。

对「相」（phase）一词的使用要谨慎，原因有二：其一是 FI 参数在重整化群下跑动，其二是场空间的不同区域还存在其他真空，例如 Coulomb 真空。当我们说「相」时，指的是这样一段能标范围中的理论：在其中 $\zeta$ 处于某个腔室内，且场空间位于稳定子群有限、经典分析成立的区域。需要说得更清楚时，我们有时会改用「区制」（regime）一词。在几何区制中，$\zeta$ 的跑动对应于 Kähler 类的跑动，$\theta$ 的轴向平移对应于 B-场的平移。

\subsection{例子}

\para{CP$^{N-1}$ 模型}

考虑如下模型：
\begin{align*}
    G &= U(1),\\
    V &= \C(1)^{\oplus N}\ni\phi=(\phi_1,\ldots,\phi_N),\\
% [原文 p.24]
    W &= 0,\\
    \tW &= -t\sigma.
\end{align*}
其中 $\C(i)$ 是 $U(1)$ 的权为 $i$ 的表示（$g\in U(1)$ 的作用是乘以 $g^i$），$t=\zeta-\ii\theta\in\C/2\pi\ii\Z$。由于 $\phi_i$ 之间不存在规范不变的多项式，不允许有超势。存在具有荷整性的矢量 $U(1)$ R-对称性，例如令所有 $\phi_i$ 的 R-荷为零。

标量势为
\begin{equation}
    U = |\sigma\phi|^2 + \frac{e^2}{2}\left(|\phi|^2-\zeta\right)^2. \tag{4.23}
\end{equation}
当 $\zeta>0$ 时，真空方程 $U=0$ 要求 $|\phi|^2=\zeta$。特别地，$\phi$ 取非零值，规范群的稳定子群是平凡的。在这种情形下，我们说「真空把规范群完全破缺掉」。另一个方程 $|\sigma\phi|^2=0$ 则要求 $\sigma$ 必须为零。事实上，这与「规范群被破缺到一个有限子群」这一事实相联系。因此，经典超对称真空的空间是满足 $|\phi|^2=\zeta$ 的 $\phi$ 的空间模去规范群作用，即复射影空间 $\CP^{N-1}$。所有横截于该真空流形的模式都具有 $e\sqrt{\zeta}$ 量级的质量，于是理论在低于 $e\sqrt{\zeta}$ 的能量下约化为靶空间为 $\CP^{N-1}$ 的非线性 sigma 模型。从拉氏量可以读出 Kähler 类与 B-场类为
\begin{equation}
    [\omega]\simeq\zeta H\in\mathrm{H}^2(\CP^{N-1},\R),\qquad
    [B]\simeq[\theta H]\in\mathrm{H}^2(\CP^{N-1},\R/2\pi\Z), \tag{4.24}
\end{equation}
其中 $H$ 是 $\CP^{N-1}$ 的超平面类。

% [Tutor Note: p.24 式 (4.24) 第二式的群系数是 $\R/2\pi\Z$ 而非 $\R$：B-场以 $2\pi$ 为周期（$\int_{\Sigma_2}B$ 只定义到 $2\pi\Z$），故 $\theta$ 的周期为 $2\pi$。核对办法：$\theta$ 只通过 $t=\zeta-\ii\theta$ 进入理论，而 $t$ 定义在 $\C/2\pi\ii\Z$ 上（$t\sim t+2\pi\ii$ 对应 $\theta\sim\theta-2\pi$）；再令 $\theta\to\theta+2\pi$，则 $[B]=[\theta H]$ 平移 $2\pi H$，正是 B-场的整周期位移（大规范变换），故 $[B]$ 落在 $\mathrm{H}^2(\CP^{N-1},\R/2\pi\Z)$ 中。]
当 $\zeta<0$ 时，方程 $U=0$ 没有解。也就是说，不存在经典超对称真空，我们说「超对称（在经典层面）破缺」。势在 $\phi=0$ 而 $\sigma$ 任意处取极小，真空能量为 $e^2\zeta^2/2$。

上述经典图像会收到显著的量子修正。首先，FI 参数按重整化群跑动为 $\zeta'=\zeta+N\log(\mu'/\mu)$。它在高能下是大正数，在低能下是大负数。这意味着，把理论解释为 $\CP^{N-1}$ sigma 模型，只在能量足够高（但低于 $e\sqrt{\zeta}$）、即 $\zeta$ 为大正数时才成立。在 $\zeta$ 趋于零的更低能量下它不再成立。在更低的能量下 $\zeta$ 变为负值。此时，如上所见，经典真空（$\phi=0$、$\sigma$ 任意）具有正能量 $e^2\zeta^2/2$，超对称破缺。然而，标量势会收到量子修正。在 $\sigma$ 很大、$\phi_i$ 多重态很重的区域，我们应当把有质量的
% [原文 p.25]
$\phi_i$ 多重态积掉，而不是简单地把它们设为零，并转而研究规范多重态的有效拉氏量。规范耦合常数变成 $\sigma$ 的函数 $e_{\mathrm{eff}}(\sigma)$，它在 $\sigma$ 很大时趋于经典值 $e$。有效扭超势可以精确算出，其结果为
\begin{equation}
    \tW_{\mathrm{eff}}(\sigma) = -t\sigma - N\sigma\log(\sigma/\mu) = -N\sigma\log(\sigma/\Lambda), \tag{4.25}
\end{equation}
其中
\begin{equation}
    \Lambda^N := \ee^{-t}\mu^N \tag{4.26}
\end{equation}
是重整化群不变的标度参数。有效势为
\begin{equation}
    U_{\mathrm{eff}}(\sigma) = \min_{n\in\Z}\frac{e_{\mathrm{eff}}^2(\sigma)}{2}
    \left|\partial_\sigma\tW_{\mathrm{eff}}(\sigma)+2\pi\ii n\right|^2, \tag{4.27}
\end{equation}
而真空方程是 $\partial_\sigma\tW_{\mathrm{eff}}(\sigma)=-N\log(\sigma/\Lambda)\equiv 0 \bmod 2\pi\ii\Z$，即
\begin{equation}
    \sigma^N = \Lambda^N. \tag{4.28}
\end{equation}
% [Tutor Note: p.25 式 (4.25)--(4.26) 的等价性可逐项核对：把 $\Lambda^N=\ee^{-t}\mu^N$ 代入 $-N\sigma\log(\sigma/\Lambda)$，得 $-N\sigma\log(\sigma/\mu)-N\sigma\log(\mu/\Lambda)$，而 $N\log(\mu/\Lambda)=N\log\mu-\log(\Lambda^N)=N\log\mu-(-t+N\log\mu)=t$，于是回到 $-t\sigma-N\sigma\log(\sigma/\mu)$。这同时说明 (4.26) 定义的 $\Lambda$ 与 $\mu$ 无关（RG 不变量）。由 (4.28) 得 $N$ 个真空 $\sigma=\Lambda\,\ee^{2\pi\ii k/N}$（$k=0,\dots,N-1$），与紧随其后「$\CP^{N-1}$ 的上同调群秩为 $N$」的陈述互相印证。]
我们看到有 $N$ 个解。也就是说，量子理论中有 $N$ 个超对称真空，这与经典分析的结论相反。在每一个解附近展开势都会出现二次项，因此每个真空只有有质量的激发——我们说「每个真空都有质量间隙」。

这是规范线性 sigma 模型的结论。然而，由于在特定的能标窗口内（只要 $e\gg\Lambda$）该模型可以解释为 $\CP^{N-1}$ sigma 模型，这也可以看作 $\CP^{N-1}$ sigma 模型的结论。确实，已知 sigma 模型的超对称真空与靶空间的上同调类一一对应，\footnote{这来自如下事实：sigma 模型可做 A-扭转，故超对称真空与 $\cR_A$ 的元素之间存在一一对应；而 $\cR_A$ 作为向量空间与靶空间的上同调群同构。}因此「有 $N$ 个超对称真空」这一结论，与「$\CP^{N-1}$ 的上同调群秩为 $N$」这一事实是相符的。
